\section{The Fitting Reconstruction}
\label{sec:fitting}

\statusestablished

The Scott and Anderson candidates keep positivity on the abstract property domain of Section~\ref{sec:property-domain}. The Fitting variant does not: it gives properties internal structure by separating intensions from their extensions and applies positivity to the latter. This section develops that third architecture, which is where the four-valued setting departs most sharply from the classical theory, because the choice of extensional domain becomes a substantive semantic commitment rather than a notational one.

\subsection{From a type distinction to a domain obstruction}
\label{sec:fitting-obstruction}

The Fitting comparison begins by distinguishing bilateral extensions from intensions:
\[
  Extension(Entity)
  \qquad\text{and}\qquad
  Intension(World,Entity)=World\to Extension(Entity).
\]
The positivity predicate is typed over extensions, and \(extensionAt\) makes the passage from an intension to its extension at a world explicit. Signed necessary entailment, essence, and necessary existence are then reconstructed over rigid extensions, reflecting the extensional architecture of the mechanized Fitting variant \cite{BenzmuellerFuenmayor2019}.

A naive first lift quantified Fitting Godlikeness and the relevant consistency/completeness package over every inhabitant of the bilateral extension type. Lean exposes a problem with that choice. The unrestricted type contains the extension \(B\) satisfying
\[
  \pos{B(x)}\quad\text{and}\quad\neginfo{B(x)}
\]
for every individual. Consequently Lean proves
\[
  CONS_{G,F}^{\mathrm{all}}
  \Longrightarrow
  \forall w\forall x\,\neg G_F^+(x,w),
\]
and therefore also
\[
  REG_{G,F}^{\mathrm{all}}
  \Longrightarrow
  \forall w\forall x\,\neg G_F^+(x,w).
\]
The earlier unrestricted implication from this regularity package to extensional essence remains a valid implication, but its intended Godlike antecedent is excluded by the package itself. I therefore do not count that unrestricted route as the substantive Fitting recovery result. The failed modeling choice is retained in the formal development because the obstruction is itself informative.

\subsection{Admissible extensions}
\label{sec:fitting-admissible}

The corrected candidate uses a selected domain of admissible extensions, written \(Adm(X)\), with FDE-negation closure
\[
  Adm(X)\Longrightarrow Adm(\neg X).
\]
This is a property-domain restriction rather than a global consistency postulate. In particular, admissible extensions may remain glutty, as the finite S5 witness above demonstrates.

Positive Fitting Godlikeness becomes
\[
  G_F^+(x,w)
  \quad\text{iff}\quad
  \forall X\bigl(Adm(X)\land \pos{P(X)}@w
  \Rightarrow \pos{X(x)}\bigr).
\]
The relevant Fitting regularity package is likewise restricted to admissible extensions. Lean then proves the non-vacuous extensional essence theorem
\[
\begin{split}
  &G\text{-admissible}
  +G_F\text{-realization}
  +A1_{L,F}^{adm}
  +REG_{G,F}^{adm}\\
  &\qquad\Longrightarrow
  \bigl(G_F^+(x,w)\Rightarrow
  Ess_F^+(ext_wG,x,w)\bigr).
\end{split}
\]
Unlike the current Scott-support T2 recovery theorem, no \(R^+\) premise occurs. Once local positivity of an admissible rigid extension has been recovered, its membership relation does not itself vary across accessible worlds.

The bilateral negative evidence is also audited. Under classical coherence on the selected extension domain, Lean proves the recovery schemata
\[
\begin{array}{rcl}
 Ent_F^-&\leftrightarrow&\neg Ent_F^+,\\
 G_F^-&\leftrightarrow&\neg G_F^+,\\
 Ess_F^-&\leftrightarrow&\neg Ess_F^+,\\
 BoxExists_F^-&\leftrightarrow&\neg BoxExists_F^+,\\
 NE_F^-&\leftrightarrow&\neg NE_F^+.
\end{array}
\]
These results validate the project-specific negative clauses at the classical boundary without imposing classicality on the four-valued candidate itself.

\subsection{De re and de dicto}
\label{sec:fitting-dere}

Let \(ext_wG\) denote the extension of the Godlikeness intension frozen at the evaluation world. With a distinguished \(NE\) intension realizing admissible-domain necessary existence and positive A5 on its current extension, Lean proves
\[
\begin{split}
  &\text{admissible Fitting stack}
  +\Diamond_{de\,re}^+\exists^E x\,ext_wG(x)\\
  &\qquad\Longrightarrow
  \Box_{de\,re}^+\exists^E x\,ext_wG(x).
\end{split}
\]
No reflexivity, symmetry, transitivity, seriality, or positivity-rigidity premise occurs. This is the project's four-valued counterpart of the K-level de-re behavior emphasized in the mechanized Fitting analysis \cite{BenzmuellerFuenmayor2019}.

The de-dicto claim is intentionally not identified with this theorem. It re-evaluates \(G\) at each accessible world. I isolate the required intension/extension bridge as
\[
  STAB_G:\qquad
  wRz\Rightarrow ext_wG\equiv ext_zG.
\]
Lean proves that \(STAB_G\) converts de-dicto possibility to the frozen de-re premise and converts the de-re necessary conclusion back to de-dicto necessary Godlikeness. Hence
\[
\begin{split}
  &\text{admissible Fitting stack}
  +STAB_G
  +\Diamond^+\exists^E x\,G_F(x)\\
  &\qquad\Longrightarrow
  \Box^+\exists^E x\,G_F(x),
\end{split}
\]
again without an S4 or S5 frame premise. The three-world finite witness shows that the stability condition is not merely decorative for this proof interface.

The minimized positive bridge admits a structural formulation. Lean defines \(R_{F,adm}^+\) as forward positivity rigidity restricted to admissible rigid extensions and \(R_{F,adm}^{+\leftarrow}\) as converse positivity transport along existing accessibility edges. It proves
\[
\begin{array}{rcl}
R_{F,adm}^++G_F\text{-realization}&\Longrightarrow&Refl_G^+,\\
R_{F,adm}^{+\leftarrow}+G_F\text{-realization}&\Longrightarrow&Pers_G^+,\\
R_{F,adm}^++Sym(R)&\Longrightarrow&R_{F,adm}^{+\leftarrow}.
\end{array}
\]
Consequently \(R_{F,adm}^++R_{F,adm}^{+\leftarrow}\) can replace primitive \(STAB_G^+\) in the positive de-dicto theorem without any frame condition. The finite oracle checks this bridge exhaustively over 256 two-world, one-entity configurations in the classical negation-closed admissible fragment. Among 112 candidates satisfying both transports, 32 are non-symmetric; converse positivity transport is therefore strictly weaker than symmetry in the tested family.
